\section{Conclusion}
\label{sec:conclusion}

This paper presented a gas cost analysis of L1 state commitment contract designs for
enterprise ZK validium on Avalanche Subnet-EVM. We compared three storage layouts---minimal,
rich, and events-only---and decomposed the gas budget to quantify the contribution of
Groth16 ZK proof verification relative to storage and logic operations.

Our primary contributions are:

\begin{enumerate}
\item \textbf{Gas decomposition:} ZK pairing verification dominates the gas budget at
71.9\% (205,600 of 285,756 gas for Layout A), establishing that storage layout optimization
operates within a narrow 28.1\% margin.

\item \textbf{Integrated verification requirement:} Delegated verification (cross-contract
call to a separate ZKVerifier) adds $\sim$56K gas from redundant storage and call overhead,
making it incompatible with a 300K gas target. Integrated verification is an architectural
requirement, not a preference.

\item \textbf{Optimal layout identification:} Layout A (minimal, 32 bytes/batch) achieves
285,756 gas per submission while maintaining on-chain root history queryability and full
invariant enforcement. It uses 33--50\% less per-batch storage than production ZK rollup
systems.

\item \textbf{Single-phase atomic commitment:} The enterprise validium trust model enables
atomic proof-and-commit in a single L1 transaction, eliminating the 2--3 phase separation
used by public rollups and reducing per-batch L1 transaction overhead.

\item \textbf{Invariant completeness:} A single root chain continuity check
($\mathit{prevRoot} = \mathit{currentRoot}$) is sufficient to enforce gap detection,
reversal prevention, and sequential ordering, confirmed by 7/7 invariant tests across
all layouts.
\end{enumerate}

The recommended production architecture uses Layout A with integrated Groth16 verification,
per-enterprise state isolation via address-keyed mappings, and batch metadata emitted
exclusively through events. This design leaves a 14,244 gas margin (4.7\%) below the
300K target, sufficient for access control modifiers and minor safety checks.

Future work should investigate recursive proof aggregation to amortize the 205,600 gas
verification cost across multiple batches, cross-enterprise state verification protocols
that leverage the on-chain root history, and migration to proof systems (PLONK, Halo2) with
lower per-verification gas costs.
